% Capitulo: Master data: de projeto eterno a oportunidade
\chapter{Master data: de projeto eterno a oportunidade}
\label{cap:master-data}

\rascunho{
  \textbf{Ideia do capítulo.} Master data (os cadastros mestres de clientes, produtos,
  fornecedores, locais) é a base de quase todo o resto do livro, e por muito tempo foi
  visto como um projeto caro, lento e difícil de justificar. Em muitas empresas era mais
  fácil continuar com as planilhas. A tese do capítulo: hoje dá para tratar master data
  como oportunidade, puxando cada domínio por um projeto de negócio que já precisa dele,
  já aprovado e com recurso. Essa é a abordagem que funcionou na minha experiência.
  \begin{itemize}
    \item Ligações: as dimensões conformadas de Kimball são master data com outro nome
      (\veja{sec:kimball}); master data costuma viver na camada prata ou ouro
      (capítulo~\ref{cap:medalhao}); e é contexto essencial para a IA saber que dois
      nomes diferentes são o mesmo cliente (capítulo~\ref{cap:contexto-ia}).
    \item Referências de base: \textcite{dama2017} e \textcite{loshin2009}; para a
      conciliação de registros, \textcite{fellegi1969} e \textcite{christen2012}.
  \end{itemize}
}

% ---------------------------------------------------------------------------
\section{O que é master data}
\label{sec:o-que-e-master-data}

\rascunho{
  \begin{itemize}
    \item Dados mestres são as entidades centrais do negócio, compartilhadas por vários
      processos e sistemas: clientes, produtos, fornecedores, locais, funcionários
      \parencite{dama2017,loshin2009}.
    \item Diferença para dado transacional (a venda, a viagem, o lançamento contábil) e
      para dado de referência (listas de códigos, como estados ou moedas).
    \item O problema concreto: o mesmo cliente aparece com nomes e IDs diferentes em cada
      sistema (Tabela~\ref{tab:registro-mestre}). Sem um registro mestre, cada relatório
      soma o cliente de um jeito.
    \item Registro mestre (\emph{golden record}) e tabela de correspondência entre os IDs
      de cada fonte e o ID mestre.
  \end{itemize}
}

\pendente{escrever}

\begin{table}[htbp]
  \centering
  \caption{O mesmo cliente em três fontes e o registro mestre que as concilia.}
  \label{tab:registro-mestre}
  \small
  \begin{tabular}{@{}lll@{}}
    \toprule
    Fonte & ID na fonte & Nome como está na fonte \\
    \midrule
    ERP                & 4471          & AGRO SILVA LTDA \\
    CRM                & C-0932        & Agro Silva \\
    Planilha comercial & (sem ID)      & Agrosilva Faz. Boa Vista \\
    \midrule
    \textbf{Registro mestre} & \textbf{1001} & \textbf{Agro Silva Ltda} \\
    \bottomrule
  \end{tabular}
  \fonte{exemplo fictício elaborado pelo autor.}
\end{table}

% ---------------------------------------------------------------------------
\section{Como master data era visto}
\label{sec:master-data-antes}

\rascunho{
  \begin{itemize}
    \item Um programa grande, separado de qualquer entrega de negócio: mapear todas as
      fontes, conciliar nomes e IDs, definir governança, e só então entregar algo.
    \item Exigia contratar analistas de dados para conversar com os analistas de negócio,
      mapear as fontes e negociar regras. Meses ou anos até ver o resultado de um master
      data bem feito (parte de cima da Figura~\ref{fig:master-data-abordagens}).
    \item Difícil de justificar: o benefício é difuso (\enquote{dados mais
      consistentes}) e o custo é concreto. A planilha de Excel com o \enquote{de-para}
      mantido à mão vencia, porque já existia e funcionava mais ou menos.
    \item O custo escondido da planilha: cada área com a sua versão, retrabalho em todo
      fechamento, números que não batem na reunião. \pendente{relato de uma situação dessas, sem nomes}
  \end{itemize}
}

\pendente{escrever}

% ---------------------------------------------------------------------------
\section{Master data puxado por um projeto de negócio}
\label{sec:master-data-puxado}

\rascunho{
  \begin{itemize}
    \item A abordagem que funcionou para mim: partir de um projeto de negócio já
      aprovado e identificar os braços dele que dependem de master data. Exemplo: um
      projeto comercial que precisa juntar a carteira de clientes de vários sistemas.
    \item Abrir um projeto de master data em paralelo, só para aquele domínio (clientes),
      com a necessidade já justificada e o recurso já garantido pelo projeto de negócio.
    \item O projeto de negócio é o primeiro cliente do master data e valida o resultado
      na prática. O próximo projeto reaproveita o domínio pronto e puxa outro (produtos,
      fornecedores). Parte de baixo da Figura~\ref{fig:master-data-abordagens}.
    \item Por que funciona: o patrocínio já existe, o escopo é pequeno, o valor aparece
      em semanas, e cada domínio entregue torna o próximo projeto mais barato.
    \item Cuidados: não deixar o prazo do projeto de negócio atropelar a qualidade do
      domínio; desenhar o domínio pensando no reuso, não só no primeiro projeto.
    \item Relato completo de um caso. \pendente{escolher o caso: o projeto de negócio, o domínio, quanto tempo levou, o que foi reaproveitado depois}
    \item Ligação com a Parte~\ref{parte:agil}: o projeto paralelo entra no mesmo quadro,
      como épico próprio, com as histórias amarradas às do projeto de negócio.
  \end{itemize}
}

\pendente{escrever}

\begin{figure}[htbp]
  \centering
  % Master data: programa unico contra dominios puxados por projetos de negocio (cap. 13).
\begin{tikzpicture}[diagrama, x=1mm, y=1mm,
    titulo/.style={anchor=west, font=\footnotesize\scshape, text=cinza},
    linha/.style={anchor=west, font=\footnotesize, text=cinza}]

  % --- antes
  \node[titulo] at (0,40) {antes: um programa de master data};
  \node[caixa, anchor=west, minimum width=26mm] at (20,31) {Mapear fontes};
  \node[caixa, anchor=west, minimum width=34mm] at (46,31) {Conciliar nomes e IDs};
  \node[caixa, anchor=west, minimum width=20mm] at (80,31) {Governança};
  \node[destaquecaixa, anchor=west, minimum width=16mm] at (100,31) {Resultado};
  \draw[seta, draw=cinza] (20,24.5) -- (116,24.5);
  \node[rotulo, anchor=north] at (68,24) {meses ou anos até o primeiro resultado};

  % --- agora
  \node[titulo] at (0,12) {agora: puxado por um projeto de negócio};
  \node[linha] at (0,3) {Negócio};
  \node[linha] at (0,-9) {Master data};
  \node[caixa, anchor=west, minimum width=44mm] (pa) at (20,3) {Projeto de negócio A};
  \node[caixa, anchor=west, minimum width=46mm] (pb) at (70,3) {Projeto de negócio B};
  \node[destaquecaixa, anchor=west, minimum width=28mm] (cl) at (26,-9) {Clientes};
  \node[destaquecaixa, anchor=west, minimum width=28mm] (pr) at (80,-9) {Produtos};
  \draw[seta] (cl.north) -- (cl.north |- pa.south);
  \draw[seta] (pr.north) -- (pr.north |- pb.south);
  \draw[seta, dashed, draw=cinza] (cl.east) -- ($(pb.south west)+(4,0)$)
    node[rotulo, pos=0.45, below right] {reuso};
  \draw[seta, draw=cinza] (20,-16.5) -- (116,-16.5);
  \node[rotulo, anchor=north] at (68,-17)
    {um domínio por vez, cada um justificado por um projeto que precisa dele};
\end{tikzpicture}

  \caption{Dois jeitos de construir master data: um programa único, que entrega no
    fim, ou um domínio por vez, puxado por projetos de negócio que já precisam dele.}
  \label{fig:master-data-abordagens}
  \fonte{elaborado pelo autor.}
\end{figure}

% ---------------------------------------------------------------------------
\section{O que a IA muda na conciliação}
\label{sec:master-data-ia}

\rascunho{
  \textbf{Ideia minha para você validar com a sua experiência.}
  \begin{itemize}
    \item A conciliação de registros (\emph{record linkage}) tem teoria antiga: o
      modelo probabilístico de \textcite{fellegi1969} e as técnicas reunidas por
      \textcite{christen2012}.
    \item O que ficou mais barato com IA: sugerir pares prováveis com nomes escritos de
      jeitos diferentes, padronizar endereços e razões sociais, explicar por que dois
      registros parecem o mesmo. A pessoa revisa e decide; a IA tira o trabalho braçal.
    \item Isso encurta exatamente a etapa que fazia o master data levar meses, e reforça
      a abordagem por domínio: conciliar um domínio por vez ficou viável dentro do prazo
      de um projeto de negócio. Um motor simples de conciliação está em
      \veja{sec:motor-conciliacao}.
    \item Cuidado com privacidade: cadastro de cliente tem dado pessoal; o que vai para o
      modelo precisa ser limpo antes (ver \veja{sec:seguranca}).
  \end{itemize}
}

\pendente{escrever}

% ---------------------------------------------------------------------------
\section{Governança mínima que funciona}
\label{sec:governanca-master-data}

\rascunho{
  \begin{itemize}
    \item Um dono de negócio por domínio, que decide as regras (o que é um cliente
      ativo, quando dois cadastros são o mesmo).
    \item Regras de criação e alteração: quem pode criar um cliente novo e onde.
    \item Testes de unicidade e de correspondência rodando a cada carga (ver
      \veja{sec:tdd}).
    \item A tabela de correspondência versionada e auditável, não numa planilha.
    \item Começar pequeno: governança do tamanho do domínio entregue, crescendo junto.
    \item A governança completa, com o steward como revisor da IA, fica para o
      capítulo~\ref{cap:governanca}; como construir a aplicação, para o
      capítulo~\ref{cap:aplicacao-propria}.
  \end{itemize}
}

\pendente{escrever}

% ---------------------------------------------------------------------------
\begin{leituras}
  \leitura{dama2017}{O guia de referência de gestão de dados. O capítulo sobre dados
    mestres e de referência dá o vocabulário e as responsabilidades.}
  \leitura{loshin2009}{Um livro inteiro sobre master data, do caso de negócio à
    arquitetura e à governança.}
  \leitura{christen2012}{Para quem vai implementar a conciliação: técnicas de
    comparação, bloqueio e avaliação da qualidade dos pares encontrados.}
  \leitura{fellegi1969}{O artigo clássico que formalizou a conciliação probabilística
    de registros, ainda base de muitas ferramentas.}
\end{leituras}
